\documentclass[UTF8,fontset=fandol,a4paper,11pt]{ctexart}
\usepackage[margin=25mm,headheight=18pt]{geometry}
\usepackage{amsmath,amssymb,bm,amsthm}
\usepackage{graphicx,subcaption}
\usepackage{booktabs,tabularx,multirow,longtable}
\usepackage{siunitx}
\usepackage{xcolor,listings}
\usepackage{fancyhdr}
\usepackage{titling,enumitem}
\setlength{\droptitle}{-2em}
\setlist{itemsep=0pt,parsep=0pt,topsep=4pt}
\usepackage{csquotes}
\usepackage[backend=biber,style=numeric,sorting=none]{biblatex}
\addbibresource{references.bib}
\usepackage[hidelinks]{hyperref}
\usepackage[nameinlink,noabbrev]{cleveref}

\newcommand{\vect}[1]{\bm{#1}}
\newcommand{\R}{\mathbb{R}}
\DeclareMathOperator{\rank}{rank}
\DeclareMathOperator*{\argmin}{arg\,min}
\newtheorem{theorem}{定理}[section]
\crefname{equation}{式}{式}
\Crefname{equation}{式}{式}
\crefname{figure}{图}{图}
\Crefname{figure}{图}{图}
\crefname{table}{表}{表}
\Crefname{table}{表}{表}
\crefname{section}{节}{节}
\Crefname{section}{节}{节}
\crefname{theorem}{定理}{定理}
\Crefname{theorem}{定理}{定理}
\renewcommand{\proofname}{证明}
\pagestyle{fancy}
\fancyhf{}
\fancyhead[L]{LaTeX 排版练习}
\fancyhead[R]{结构、公式与图表}
\fancyfoot[C]{\thepage}
\renewcommand{\headrulewidth}{0.4pt}
\setlength{\parindent}{2em}
\setlength{\parskip}{0.2em}
\setcounter{tocdepth}{1}
\sisetup{detect-all}
\lstset{
  language=Python,
  basicstyle=\small\ttfamily,
  keywordstyle=\color{blue!65!black},
  commentstyle=\color{green!40!black},
  stringstyle=\color{red!60!black},
  numbers=left,
  numberstyle=\scriptsize\color{gray},
  numbersep=8pt,
  frame=single,
  rulecolor=\color{gray!50},
  breaklines=true,
  columns=fullflexible,
  keepspaces=true,
  showstringspaces=false,
  tabsize=4
}
\hypersetup{
  pdftitle={LaTeX 排版：可复现练习工程},
  pdfauthor={LSQY},
  pdfsubject={中文文档、数学公式、图表与参考文献}
}
\title{\LaTeX{} 排版\\[0.4em]\large 可复现练习工程}
\author{LSQY 学习笔记}
\date{2026 年 10 月 3 日}

\begin{document}
\maketitle
\thispagestyle{fancy}
\begin{abstract}
本工程将一份中文技术文档拆成结构、数学、图表和参考文献四个部分。
公式、图片与书目都有可编辑的源文件；通过修改内容观察自动编号与版面的变化，
比记忆孤立命令更容易掌握排版规则。这里的机器人关节数据为教学构造数据，
不代表实际测量或控制器性能。
\end{abstract}
\tableofcontents
\section{从结构开始}\label{sec:structure}
正文只描述内容的层级与含义。章、节、图、表、公式的序号由计数器管理，
引用目标由标签标识。不要用手打空格对齐正文，也不要把编号写死在句子中。
例如数学部分见\cref{sec:math}，坐标图见\cref{fig:coordinate-pair}。

\begin{itemize}
  \item \texttt{main.tex}：文档类、宏包、字体、页眉与章节装配。
  \item \texttt{sections/}：按学习主题组织的正文源文件。
  \item \texttt{figures/}：坐标图的 TikZ 源文件及编译生成的 PDF。
  \item \texttt{references.bib}：真实文献的元数据，引用键不依赖显示顺序。
\end{itemize}

\textbf{第一项练习：}在本节前增加一个小节，重新构建后检查目录、正文引用与页码。
一次构建应让依赖稳定；PDF 能打开，只能说明产物存在，不能替代检查警告和版面。
学习排版底层模型可以参阅 Knuth 的著作\cite{knuth1984texbook}。

\clearpage
\section{最小二乘与数学结构}\label{sec:math}
设 \(\vect{A}\in\R^{m\times n}\)、\(\vect{x}\in\R^n\)、
\(\vect{b}\in\R^m\)，所有元素均为实数。最小二乘目标为
\begin{equation}
  J(\vect{x})=\frac{1}{2}\lVert\vect{A}\vect{x}-\vect{b}\rVert_2^2.
  \label{eq:least-squares}
\end{equation}
两条分别编号的公式使用 \texttt{align}，以等号为对齐点：
\begin{align}
  \nabla J(\vect{x})&=\vect{A}^{\mathsf{T}}
    (\vect{A}\vect{x}-\vect{b}),\label{eq:gradient}\\
  \nabla^2J(\vect{x})&=\vect{A}^{\mathsf{T}}\vect{A}.
    \label{eq:hessian}
\end{align}
一条推导需要多行、但只保留一个编号时，使用外层公式包住 \texttt{aligned}：
\begin{equation}
  \begin{aligned}
    J(\vect{x})
      &=\frac{1}{2}(\vect{A}\vect{x}-\vect{b})^{\mathsf{T}}
        (\vect{A}\vect{x}-\vect{b})\\
      &=\frac{1}{2}\vect{x}^{\mathsf{T}}\vect{A}^{\mathsf{T}}
        \vect{A}\vect{x}
        -\vect{b}^{\mathsf{T}}\vect{A}\vect{x}
        +\frac{1}{2}\vect{b}^{\mathsf{T}}\vect{b}.
  \end{aligned}
  \label{eq:expanded-loss}
\end{equation}
以阈值 \(\delta>0\) 定义分段损失，条件列与表达式列分开对齐：
\begin{equation}
  \rho_\delta(r)=
  \begin{cases}
    \dfrac{1}{2}r^2,&\lvert r\rvert\leq\delta,\\
    \delta\bigl(\lvert r\rvert-\dfrac{1}{2}\delta\bigr),
      &\lvert r\rvert>\delta.
  \end{cases}
  \label{eq:huber}
\end{equation}
矩阵环境中的分列、换行同样表达结构，例如
\[
  \vect{A}=\begin{bmatrix}1&0\\1&1\\1&2\end{bmatrix},
  \qquad
  \vect{x}=\begin{bmatrix}a\\b\end{bmatrix},
  \qquad
  \vect{b}=\begin{bmatrix}1\\2\\2\end{bmatrix}.
\]
\begin{theorem}[正定二次型]\label{thm:positive-quadratic}
若实对称矩阵 \(\vect{Q}\) 正定，则
\(f(\vect{x})=\vect{x}^{\mathsf{T}}\vect{Q}\vect{x}\)
的唯一全局最小点为 \(\vect{x}=\vect{0}\)。
\end{theorem}
\begin{proof}
由正定性，非零向量满足 \(f(\vect{x})>0\)，而 \(f(\vect{0})=0\)。
\end{proof}
练习：引用\cref{eq:gradient,eq:hessian}，再插入新公式观察编号变化。
将\cref{thm:positive-quadratic}的“正定”改成“半正定”后，数学结论需要重新检查。

\clearpage
\section{图片、关节数据与分组表格}\label{sec:figures-tables}
图题在图后，标签放在图题之后。这里重复使用同一份矢量图，
演示两个独立子图的编号与引用；它们不是不同实验结果。
\begin{figure}[htbp]
  \centering
  \begin{subfigure}[b]{0.44\linewidth}
    \centering
    \includegraphics[width=0.88\linewidth]{figures/coordinate-frame.pdf}
    \caption{坐标系示意}\label{fig:frame-a}
  \end{subfigure}
  \hfill
  \begin{subfigure}[b]{0.44\linewidth}
    \centering
    \includegraphics[width=0.88\linewidth]{figures/coordinate-frame.pdf}
    \caption{重复引用同一源图}\label{fig:frame-b}
  \end{subfigure}
  \caption{一份矢量资源组成的双子图练习}\label{fig:coordinate-pair}
\end{figure}

表格中的数字按小数点对齐，单位统一放在表头。下面使用
\texttt{tabularx} 的弹性文字列和 \texttt{siunitx} 的数字列。
\begin{table}[htbp]
  \centering
  \caption{教学用关节跟踪误差，数据为人为构造}\label{tab:joint-error}
  \begin{tabularx}{\linewidth}{@{}l X S[table-format=1.3]
    S[table-format=1.3]@{}}
    \toprule
    关节 & 说明 & {RMSE (\si{\radian})} & {峰值 (\si{\radian})}\\
    \midrule
    J1 & 髋部俯仰，仅用于演示排版 & 0.012 & 0.031\\
    J2 & 膝部俯仰，数值保留三位小数 & 0.008 & 0.024\\
    J3 & 踝部俯仰，单位与前两项一致 & 0.015 & 0.040\\
    \bottomrule
  \end{tabularx}
\end{table}

跨行表示一个组包含多个条目，跨列表示共同的表头。三线表通常不加竖线，
组间局部横线用 \texttt{cmidrule} 表达，避免让表格过于密集。
\begin{table}[htbp]
  \centering
  \caption{跨行与跨列结构示例}\label{tab:grouped}
  \begin{tabular}{@{}llrr@{}}
    \toprule
    模块 & 项目 & \multicolumn{2}{c}{样本数量}\\
    \cmidrule(l){3-4}
         &      & 训练 & 检查\\
    \midrule
    \multirow{2}{*}{关节} & 位置 & 120 & 30\\
                         & 速度 & 120 & 30\\
    \midrule
    \multirow{2}{*}{姿态} & 横滚 & 80 & 20\\
                         & 俯仰 & 80 & 20\\
    \bottomrule
  \end{tabular}
\end{table}

练习：将\cref{tab:joint-error}的说明改成较长句子，观察文字列换行；
把某个数字改为两位整数时，同时调整 \texttt{table-format}。
浮动体的 \texttt{htbp} 是位置偏好，不能保证恰好放在源码所在行。

\clearpage
\section{长表格、代码与验证}\label{sec:long-document}
跨页表格直接使用 \texttt{longtable}，不要再套 \texttt{table} 浮动环境。
本例先保留少量行；复制数据行后，继续页会显示重复表头。

\begin{longtable}{@{}p{0.12\linewidth}p{0.34\linewidth}p{0.45\linewidth}@{}}
  \caption{练习记录：包含继续页的表头与页尾定义}\label{tab:practice-log}\\
  \toprule
  序号 & 修改内容 & 验收现象\\
  \midrule
  \endfirsthead
  \multicolumn{3}{l}{表 \thetable：练习记录（续）}\\
  \toprule
  序号 & 修改内容 & 验收现象\\
  \midrule
  \endhead
  \midrule
  \multicolumn{3}{r}{续下页}\\
  \endfoot
  \bottomrule
  \endlastfoot
  01 & 新增带编号公式 & 正文引用自动更新\\
  02 & 调整图的宽度 & 图未越过正文边界\\
  03 & 加长表格说明 & 文字列自动换行\\
  04 & 增加真实书目 & 引用键解析为文献编号\\
  05 & 复制本表的数据行 & 跨页后显示继续页表头\\
\end{longtable}

下面的 Python 代码只含 ASCII 字符，避免把中文字体问题与代码环境问题混在一起。
它实现\cref{eq:huber}；命令的语法高亮不能证明算法逻辑正确。
\begin{lstlisting}[caption={Huber loss with an explicit threshold},label={lst:huber}]
def huber(residual, delta=1.0):
    if delta <= 0:
        raise ValueError("delta must be positive")
    magnitude = abs(residual)
    if magnitude <= delta:
        return 0.5 * residual * residual
    return delta * (magnitude - 0.5 * delta)

assert huber(0.0) == 0.0
assert huber(1.0) == 0.5
assert huber(-2.0) == 1.5
\end{lstlisting}

\textbf{一次完整检查：}
确认没有未定义引用、缺失字符或越界盒子；逐页观察公式、表格和图题；
点击一个公式引用和一条文献引用；最后核对数学条件、单位及数据来源。
工程构建使用 \texttt{latexmk} 驱动 XeLaTeX 与 Biber，
而不是把“生成了 PDF”作为唯一成功条件。

\clearpage
\section{引用链路与下一次练习}\label{sec:references}
引用键是文献的稳定标识，显示出来的数字由引用顺序决定。
本例先引用 Knuth\cite{knuth1984texbook}，
再引用 Lamport 对文档准备系统的介绍\cite{lamport1994latex}。
两条记录都对应真实出版物；示例中的机器人数据与这些书目没有实验来源关系。

本工程采用 BibLaTeX 加 Biber：
\begin{enumerate}
  \item XeLaTeX 写出文献请求与交叉引用辅助信息。
  \item Biber 读取 \texttt{references.bib}，生成处理后的文献数据。
  \item XeLaTeX 读入数据，继续编译直到引用稳定。
\end{enumerate}
平时用 \texttt{latexmk -xelatex main.tex} 让工具按依赖判断构建步骤。
如果出现未定义引文，应检查引用键、文件路径、Biber 日志和版本配套，
不能靠手写文献编号补救。

\subsection*{交付前的四个问题}
\begin{itemize}
  \item 结构：新增章节后，目录、标签和超链接是否仍正确？
  \item 排版：长公式、数字列、续表和页眉是否均在预期位置？
  \item 内容：符号维度、定理条件、数据单位和参考来源是否准确？
  \item 复现：换到同一构建环境，是否能从源文件重新生成 PDF？
\end{itemize}

\printbibliography[heading=subbibliography,title={参考文献}]

\bigskip
\noindent
建议每次只改变一种因素，并保留修改前后的 PDF。
本练习中的显式换页用于将每个主题独立展示；
正式长文档应先让结构与浮动体自然排版，再针对需要处理分页。

\end{document}
